\documentclass[11pt,a4paper]{article}
\usepackage[utf8]{inputenc}
\usepackage[T1]{fontenc}
\usepackage[brazil]{babel}
\usepackage[a4paper,margin=2.3cm]{geometry}
\usepackage{enumitem}
\usepackage{hyperref}
\usepackage{xcolor}
\setlength{\parindent}{0pt}
\setlength{\parskip}{5pt}
\sloppy
\setlength{\emergencystretch}{3em}
\hypersetup{colorlinks=true,urlcolor=blue,linkcolor=black}
\newcommand{\cabecalho}[2]{\begin{center}{\small \textbf{Sprint 2 -- #1}}\\[1mm]\textit{Inventário de Ativos e Vulnerabilidades}\\Matheus Moreli Durão Silva -- 12621CBS217\\Cibersegurança -- Bacharelado -- Noturno\end{center}\hrule\vspace{4mm}}
\newcommand{\bloco}[1]{\par\vspace{2mm}\textbf{#1}\par}
\newcommand{\codigo}[1]{\texttt{#1}}

\begin{document}\cabecalho{S2\_01 -- Boas Práticas de Código}{}
\bloco{O que o card pede} Código legível, nomes que expressem intenção, funções com responsabilidade principal, sem segredos, tratamento seguro de erro e preservação do comportamento anterior.
\bloco{Como funciona} Uma boa prática não é uma ``função secreta''. É uma escolha que permite ler o programa em ordem: \\ \codigo{cadastrar\_ativo} cria; \codigo{consultar\_por\_id} busca; \codigo{salvar\_ativos} grava. O nome substitui a necessidade de adivinhar o que o trecho faz.
\bloco{Aplicação no projeto} O código foi dividido em \codigo{app.py}, \codigo{interface.py}, \codigo{inventario.py}, \codigo{modelos.py} e \codigo{persistencia.py}. A tela usa \codigo{input} e \codigo{print}; as regras ficam em \codigo{inventario.py}; a gravação JSON fica isolada. Não há senha, token ou credencial no código. Erros conhecidos retornam mensagens curtas, sem exibir detalhes internos.
\bloco{Como conferir} Ler uma função por vez e perguntar: ela faz uma tarefa principal? Executar \codigo{python -m unittest -v test\_sprint2.py}; os testes confirmam que o CRUD anterior continua funcionando após a divisão.
\bloco{Checklist atendido} Nomes claros; responsabilidade única; nenhum dado sensível; mensagens seguras; testes de regressão executados.
\end{document}
